\par\SourceLeaf{525}Je seltener diese Neigung ist, desto mehr freut es mich, dass Sie sich dem Teil der Mathematik so stark verbunden fühlen, der seit jeher mein Lieblingsstudium ist. Ich wünsche Ihnen von Herzen äußere Verhältnisse, unter denen Sie möglichst weitgehend Herr Ihrer Zeit und der Wahl Ihrer Arbeiten bleiben. Mich selbst hinderten gleich nach dem Erscheinen meiner Disquisitiones andersartige Beschäftigungen und später meine äußeren Verhältnisse sehr daran, meiner Neigung in dem gewünschten Maß zu folgen. Statt eines zweiten Teils dieses Werks, den ich früher beabsichtigte, werde ich mich aller Wahrscheinlichkeit nach darauf beschränken müssen, von Zeit zu Zeit eine Abhandlung über einen einzelnen Gegenstand vorzulegen.
Die drei bisher erschienenen Abhandlungen dieser Art im 16. Band der hiesigen Commentationen und im ersten und vierten Band der Commentationes recentiores enthalten jedoch (abgesehen von einem Teil der zweiten) keine der Fragen, die ich schon 1801 für die Fortsetzung im Blick hatte, sondern neue. Auch meine späteren Arbeiten dieser Art betreffen daher einen neuen Gegenstand, nämlich die Theorie der biquadratischen Reste, die ich in etwa drei Abhandlungen darzustellen denke. Die erste davon wird bald für den sechsten Band der Comment. rec. gedruckt werden; die Hauptmaterialien für das Übrige sowie für die ähnliche Theorie der kubischen Reste sind im Wesentlichen als erledigt anzusehen, obwohl davon noch wenig ordentlich zu Papier gebracht ist.
\par Empfehlen Sie mich bitte Herrn von \textsc{Humboldt}, falls er noch in Paris ist, und entschuldigen Sie, dass ich ihm jetzt nicht selbst schreibe, mit meiner Sorge, mein Brief könnte ihn nicht erreichen, da er, wie ich höre, Paris verlassen wollte.
\par\Needspace{5\baselineskip}Mit aufrichtiger Hochschätzung
\par\nobreak\noindent\begin{tabular*}{\linewidth}{@{}l@{\extracolsep{\fill}}r@{}}
&Ihr ergebenster\\
Göttingen, den 13. September 1826.&C. F. \textsc{Gauss}.\\
\end{tabular*}
\BookRule{.25\textwidth}
\par\endgroup
